\documentclass[main]{subfiles}

\begin{document}
\ifSubfilesClassLoaded{\setcounter{chapter}{1}}{}
\chapter{Theoretical Foundations}

\section{Overview}

Chapter~1 introduced, at a survey level, the two theoretical frameworks this
thesis unifies -- the Standard Model Extension (SME) and the
Dobrescu--Mocioiu (DM) classification of exotic potentials -- along with the
Foldy--Wouthuysen (FW) transformation that bridges them. This chapter
develops each of the three in full: Section~\ref{sec:sme} presents the
general minimal SME fermion Lagrangian and classifies all of its
coefficients by CPT and Lorentz behaviour, not only the three ($b_\mu$,
$H_{\mu\nu}$, $d_{\mu\nu}$) this thesis derives explicitly. Section~\ref{sec:dm}
gives the complete, closed-form $V_1$--$V_{16}$ DM potential catalogue used
throughout the rest of the thesis. Section~\ref{sec:fw} develops the general
FW transformation machinery -- the generator, the order-by-order expansion,
and the even/odd operator decomposition -- independent of any particular SME
coefficient. Chapter~3 then applies this machinery to $b_\mu$, $H_{\mu\nu}$,
and $d_{\mu\nu}$ specifically, deriving the explicit SME$\to$DM mapping that
is the central technical contribution of this thesis.

\section{The Standard Model Extension}
\label{sec:sme}

\subsection{Effective Field Theory Framework}

The SME is constructed as the most general effective field theory
Lagrangian obtainable by adding to the Standard Model (and General
Relativity) all observer-scalar terms formed by contracting Standard Model
field operators with coefficients carrying free Lorentz indices
\cite{Colladay1997,Colladay1998}. A subtlety central to the entire
framework is the distinction between two kinds of Lorentz transformation.
\emph{Observer} Lorentz transformations rotate or boost the coordinate frame
used to describe a physical system, without touching the system itself; the
SME is constructed to respect observer Lorentz invariance exactly; every term
in the SME Lagrangian is a genuine Lorentz scalar under this class of
transformation; if it were not, the theory would depend on the physicist's
choice of coordinates, which is never observed. \emph{Particle} Lorentz
transformations, by contrast, rotate or boost a localised particle or field
configuration while holding the background coefficients (and hence the
preferred directions or frames they define) fixed. It is only under particle
Lorentz transformations that the SME's background coefficients -- $a_\mu$,
$b_\mu$, $c_{\mu\nu}$, $d_{\mu\nu}$, $H_{\mu\nu}$, and so on, in the fermion
sector -- act as fixed, non-dynamical tensors that break the symmetry. This
distinction is what allows the SME to be a coordinate-independent, internally
consistent field theory while still describing genuine, in-principle
observable violations of the Lorentz invariance normally assumed by the
Standard Model and General Relativity.

The SME Lagrangian is written schematically as
\begin{equation}
  \mathcal{L}_{\mathrm{SME}} = \mathcal{L}_{\mathrm{SM}} + \delta\mathcal{L}_{\mathrm{LV}},
\end{equation}
where $\mathcal L_{\mathrm{SM}}$ is the ordinary Standard Model Lagrangian
and $\delta\mathcal L_{\mathrm{LV}}$ collects every observer-Lorentz-scalar,
particle-Lorentz-violating operator that can be built from Standard Model
fields, organised by mass dimension. This thesis works exclusively within
the \emph{minimal} SME: the subset of $\delta\mathcal L_{\mathrm{LV}}$
restricted to operators of renormalisable mass dimension (dimension $\leq4$)
in the fermion sector, which is where essentially all existing
spin-dependent exotic-interaction experiments place their bounds.

\subsection{The Minimal SME Fermion Lagrangian}

For a single Dirac fermion $\psi$ of mass $m$, the general minimal SME
fermion-sector Lagrangian is
\begin{equation}
  \mathcal L_{\mathrm{fermion}}
  = \frac{i}{2}\bar\psi\,\Gamma^\nu\overleftrightarrow{D}_\nu\,\psi
  - \bar\psi\,M\,\psi, \label{eq:sme-general}
\end{equation}
where $D_\nu$ is the gauge-covariant derivative,
$\overleftrightarrow{D}_\nu \equiv D_\nu - \overleftarrow{D}_\nu$, and the
Lorentz-violating physics is packaged entirely into two composite
matrix-valued objects,
\begin{align}
  \Gamma^\nu &= \gamma^\nu + c^{\mu\nu}\gamma_\mu + d^{\mu\nu}\gamma_5\gamma_\mu
               + e^\nu + if^\nu\gamma_5 + \tfrac12 g^{\lambda\mu\nu}\sigma_{\lambda\mu},
               \label{eq:Gamma-general} \\
  M &= m + a_\mu\gamma^\mu + b_\mu\gamma_5\gamma^\mu + \tfrac12 H_{\mu\nu}\sigma^{\mu\nu}.
      \label{eq:M-general}
\end{align}
$\Gamma^\nu$ modifies the kinetic (derivative-coupling) sector, while $M$
modifies the mass sector. Every coefficient in Eqs.~\eqref{eq:Gamma-general}
and \eqref{eq:M-general} is a fixed background tensor, not a dynamical
field, and each is in principle independent for every Standard Model fermion
species. Setting all coefficients to zero except $m$ recovers the ordinary
Dirac Lagrangian. For the leading-order (dimension $\leq4$) subset used
throughout this thesis, Eq.~\eqref{eq:sme-general} reduces to the compact
form already introduced in Chapter~1,
\begin{multline}
  \mathcal L_{\mathrm{fermion}}
  = \bar\psi\left(i\gamma^\mu D_\mu - m\right)\psi
  - a_\mu\,\bar\psi\gamma^\mu\psi
  - b_\mu\,\bar\psi\gamma_5\gamma^\mu\psi \\
  - \tfrac12 H_{\mu\nu}\,\bar\psi\sigma^{\mu\nu}\psi
  - c_{\mu\nu}\,\bar\psi\gamma^\mu D^\nu\psi
  - d_{\mu\nu}\,\bar\psi\gamma_5\gamma^\mu D^\nu\psi,
  \label{eq:sme-minimal}
\end{multline}
which is the version this thesis derives from directly. Note the operator
ordering in the $b_\mu$ term: the coefficient multiplies
$\gamma_5\gamma^\mu$, not $\gamma^\mu\gamma_5$. Since $\{\gamma_5,\gamma^\mu\}=0$,
these two orderings differ by an overall sign, so this convention choice is
not cosmetic -- it fixes the sign of every non-relativistic result derived
from $b_\mu$ in Chapter~3, and is used consistently throughout this thesis.

\subsection{Classification of SME Coefficients}

Table~\ref{tab:sme-classification} classifies every coefficient appearing in
Eq.~\eqref{eq:sme-minimal}, together with the two dimension-5 coefficients
$e_\mu,f_\mu$ that become relevant once velocity-dependent couplings are
considered, by its behaviour under the full discrete CPT transformation and
by whether it is even or odd under a particle Lorentz transformation (all of
them are Lorentz-odd, by construction, but are listed for completeness since
this is the property that makes them SME coefficients in the first place).

\begin{table}[h]
\centering
\begin{tabular}{@{}llc>{\raggedright\arraybackslash}p{6.5cm}@{}}
\toprule
Coefficient & CPT & Mass dim.\ & Leading spin-dependent role \\
\midrule
$a_\mu$        & Odd  & 4 & None at leading order; velocity-dependent at $O(1/m)$ \\
$b_\mu$        & Odd  & 4 & Zeeman-like spin coupling, $V_2$ (Ch.~3) \\
$c_{\mu\nu}$   & Even & 4 & Direction-dependent kinetic modification; subleading $V_2$-type \\
$d_{\mu\nu}$   & Even & 4 & Spin-velocity and mass-enhanced couplings, $V_2,V_7,V_8$ (Ch.~3) \\
$e_\mu$        & Odd  & 5 & Velocity-dependent, contributes to $V_8$ and higher \\
$f_\mu$        & Odd  & 5 & Velocity-dependent, contributes to $V_8$ and higher \\
$g_{\lambda\mu\nu}$ & Odd & 5 & Not treated in this thesis \\
$H_{\mu\nu}$   & Even & 4 & Anisotropic spin coupling, $V_3,V_7$ (Ch.~3) \\
\bottomrule
\end{tabular}
\caption{Classification of minimal (and near-minimal) SME fermion-sector
coefficients. CPT parity is derived case by case in Chapter~3 for $b_\mu$,
$H_{\mu\nu}$, and $d_{\mu\nu}$; the remaining entries follow the same method
and are quoted from Kostelecký and Lane (1999) and Kostelecký and Mewes
(2001).}
\label{tab:sme-classification}
\end{table}

\subsection{CPT Properties: General Argument}
\label{sec:cpt-general}

The CPT parity of a given SME coefficient is determined by how its
associated fermion bilinear transforms under the combined charge
conjugation, parity, and time-reversal operation. For a bilinear of the
schematic form $\bar\psi\,\Gamma\,\psi$, where $\Gamma$ is some product of
Dirac matrices, the standard bilinear covariants transform under CPT as
\begin{equation}
  \mathrm{CPT}\!\left[\bar\psi\,\Gamma\,\psi\right]
  = \eta_\Gamma\,\bar\psi\,\Gamma\,\psi,
  \qquad
  \eta_\Gamma = \begin{cases}
    +1 & \Gamma \in \{\mathbb 1,\; \gamma^\mu,\; \sigma^{\mu\nu}\} \\
    -1 & \Gamma \in \{\gamma_5,\; \gamma_5\gamma^\mu\}
  \end{cases}
\end{equation}
so that a coefficient multiplying a scalar, vector, or tensor bilinear
($a_\mu\bar\psi\gamma^\mu\psi$, $H_{\mu\nu}\bar\psi\sigma^{\mu\nu}\psi$) is
CPT-odd whenever the associated Lagrangian term itself must be CPT-odd to
be allowed -- and CPT-even whenever the bilinear structure is already
CPT-even, as for $c_{\mu\nu}\bar\psi\gamma^\mu D^\nu\psi$ and
$d_{\mu\nu}\bar\psi\gamma_5\gamma^\mu D^\nu\psi$. This general classification
underlies every case-by-case CPT argument made explicitly in Chapter~3, and
also underlies the general CPT rule for the non-relativistic Hamiltonians
themselves,
\begin{equation}
  H_{NR}^{\bar f}\!\left(X^{\text{CPT-odd}}\right)  = -H_{NR}^{f}\!\left(X^{\text{CPT-odd}}\right),
  \qquad
  H_{NR}^{\bar f}\!\left(X^{\text{CPT-even}}\right) = +H_{NR}^{f}\!\left(X^{\text{CPT-even}}\right),
  \label{eq:cpt-rule-general}
\end{equation}
which Chapter~3 derives explicitly for $b_\mu$ via charge conjugation and
then applies, by the general argument above, to $H_{\mu\nu}$ and
$d_{\mu\nu}$ without repeating the full charge-conjugation calculation for
each.

\section{The Dobrescu--Mocioiu Classification of Exotic Potentials}
\label{sec:dm}

\subsection{Motivation and General Structure}

Experimental searches for new long-range forces are conducted in the
laboratory, non-relativistic frame, not in the relativistic language in
which the SME (or any other beyond-Standard-Model Lagrangian) is naturally
written. Dobrescu and Mocioiu \cite{Dobrescu2006} closed this gap by
enumerating, from first principles, \emph{every} non-relativistic potential
that can arise from the exchange of a single boson of arbitrary spin and
parity between two spin-$\tfrac12$ fermions, organised by which
combinations of the two fermions' spins, relative position, and relative
velocity each potential depends on. The result is a complete basis of
sixteen potentials,
\begin{equation}
  V(\mathbf r) = \sum_{i=1}^{16} V_i(\mathbf r),
\end{equation}
each classified by its behaviour under parity ($P$) and time reversal
($T$), so that any single-boson-exchange interaction -- whatever its
relativistic origin -- can be written as some linear combination of the
$V_i$ with model-dependent coefficients $g_s,g_p,g_V,g_A$ (scalar,
pseudoscalar, vector, axial-vector). This is the language in which nearly
all spin-dependent exotic-force experiments report their bounds, which is
why matching the SME's relativistic coefficients onto this basis
(Chapter~3) is the essential step connecting theory to the experimental
constraint database compiled in Chapter~4.

\subsection{Notation and Conventions}

Table~\ref{tab:dm-notation} fixes the notation used for the potential
catalogue below and throughout the rest of this thesis.

\begin{table}[h]
\centering
\begin{tabular}{@{}l>{\raggedright\arraybackslash}p{9.5cm}@{}}
\toprule
Symbol & Meaning \\
\midrule
$\boldsymbol\sigma_{1,2}$ & Pauli matrices acting on particle 1 or 2 \\
$\hat{\mathbf r}=\mathbf r/|\mathbf r|$ & Unit vector from source to field point \\
$\lambda=1/m_\phi$ & Mediator Compton wavelength (interaction range) \\
$\mathbf v=\mathbf v_{\text{rel}}$ & Relative velocity of the two particles \\
$f,\bar f$ & Matter fermion and antimatter fermion \\
$\mathcal H_B^k=\tfrac12\varepsilon^{ijk}H_{ij}$ & Dual magnetic-like vector from $H_{\mu\nu}$ \\
$b_\mu, H_{\mu\nu}, d_{\mu\nu}$ & SME coefficients: CPT-odd vector; CPT-even tensor; CPT-even axial tensor \\
\bottomrule
\end{tabular}
\caption{Notation used throughout the Dobrescu--Mocioiu potential catalogue and Chapter~3.}
\label{tab:dm-notation}
\end{table}

\subsection{Complete Catalogue: The Sixteen Potentials}
\label{sec:dm-catalogue}

The forms below are transcribed directly from Cong et al.\ (2025), Eqs.
(1)--(12) -- equivalent to Dobrescu and Mocioiu (2006), Eqs.\
(3.1)--(3.16) -- with $y(r)=e^{-r/\lambda}/(4\pi)$, $\lambda=1/m_\phi$;
$\hat{\mathbf r}$ points from particle 2 to particle 1, $\mathbf v$ is their
relative velocity, and $m$ is the fermion mass relevant to the term. The
differential operators $(1-r\,d/dr)$ and
$(1-r\,d/dr+\tfrac13r^2d^2/dr^2)$ act on $y(r)$; evaluated explicitly they
contribute the extra $1/(\lambda r)$, $1/(\lambda^2r^2)$ terms seen in $V_3$
below (worked out in full) and analogously for $V_4$--$V_7$, $V_{15}$,
$V_{16}$ (left in operator form here; see Cong et al., Sec.~V, for the
fully evaluated closed forms once specialised to a given coupling channel).

{\small
\begin{align}
  V_1 &= \frac1r\,y(r) &&\text{(monopole--monopole, spin-independent)} \\
  V_2 &= \frac1r\,(\boldsymbol\sigma_1\cdot\boldsymbol\sigma_2)\,y(r) &&\text{(dipole--dipole / spin--spin)} \\
  V_3 &= \frac{1}{m^2r^3}\left[\boldsymbol\sigma_1\cdot\boldsymbol\sigma_2\left(1-r\frac{d}{dr}\right)
    -3(\boldsymbol\sigma_1\cdot\hat{\mathbf r})(\boldsymbol\sigma_2\cdot\hat{\mathbf r})
    \left(1-r\frac{d}{dr}+\frac{r^2}{3}\frac{d^2}{dr^2}\right)\right]y(r) \notag\\
  &= \frac{e^{-r/\lambda}}{4\pi m^2}\left\{\boldsymbol\sigma_1\cdot\boldsymbol\sigma_2
    \left[\frac1{r^3}+\frac1{\lambda r^2}\right]
    -(\boldsymbol\sigma_1\cdot\hat{\mathbf r})(\boldsymbol\sigma_2\cdot\hat{\mathbf r})
    \left[\frac3{r^3}+\frac3{\lambda r^2}+\frac1{\lambda^2r}\right]\right\}
    &&\text{(tensor dipole--dipole)} \\
  V_{4,5} &= -\frac{1}{2mr^2}(\boldsymbol\sigma_1\pm\boldsymbol\sigma_2)\cdot(\mathbf v\times\hat{\mathbf r})
    \left(1-r\frac{d}{dr}\right)y(r) &&\text{(spin--velocity; $+$=$V_4$, $-$=$V_5$)} \\
  V_{6,7} &= -\frac{1}{2mr^2}\Big[(\boldsymbol\sigma_1\cdot\mathbf v)(\boldsymbol\sigma_2\cdot\hat{\mathbf r})
    \pm(\boldsymbol\sigma_1\cdot\hat{\mathbf r})(\boldsymbol\sigma_2\cdot\mathbf v)\Big]
    \left(1-r\frac{d}{dr}\right)y(r) &&\text{($+$=$V_6$, $-$=$V_7$)} \\
  V_8 &= \frac1r(\boldsymbol\sigma_1\cdot\mathbf v)(\boldsymbol\sigma_2\cdot\mathbf v)\,y(r) &&\text{(velocity--velocity)} \\
  V_{9,10} &= -\frac{1}{2mr^2}(\boldsymbol\sigma_1\pm\boldsymbol\sigma_2)\cdot\hat{\mathbf r}
    \left(1-r\frac{d}{dr}\right)y(r) &&\text{(monopole--dipole; $+$=$V_9$, $-$=$V_{10}$)} \\
  V_{11} &= -\frac{1}{mr^2}(\boldsymbol\sigma_1\times\boldsymbol\sigma_2)\cdot\hat{\mathbf r}
    \left(1-r\frac{d}{dr}\right)y(r) &&\text{(spin cross-product; not $L\cdot S$)} \\
  V_{12,13} &= \frac{1}{2r}(\boldsymbol\sigma_1\pm\boldsymbol\sigma_2)\cdot\mathbf v\,y(r) &&\text{($+$=$V_{12}$, $-$=$V_{13}$)} \\
  V_{14} &= \frac1r(\boldsymbol\sigma_1\times\boldsymbol\sigma_2)\cdot\mathbf v\,y(r) \\
  V_{15} &= -\frac{3}{2m^2r^3}\Big\{[\boldsymbol\sigma_1\cdot(\mathbf v\times\hat{\mathbf r})](\boldsymbol\sigma_2\cdot\hat{\mathbf r}) \notag\\
    &\qquad{}+(\boldsymbol\sigma_1\cdot\hat{\mathbf r})[\boldsymbol\sigma_2\cdot(\mathbf v\times\hat{\mathbf r})]\Big\}
    \left(1-r\frac{d}{dr}+\frac{r^2}{3}\frac{d^2}{dr^2}\right)y(r) \\
  V_{16} &= -\frac{1}{2mr^2}\Big\{[\boldsymbol\sigma_1\cdot(\mathbf v\times\hat{\mathbf r})](\boldsymbol\sigma_2\cdot\mathbf v) \notag\\
    &\qquad{}+(\boldsymbol\sigma_1\cdot\mathbf v)[\boldsymbol\sigma_2\cdot(\mathbf v\times\hat{\mathbf r})]\Big\}
    \left(1-r\frac{d}{dr}\right)y(r)
\end{align}
}

Three structural notes on this catalogue are worth stating explicitly, since
each is easy to transcribe incorrectly. First, $V_2$ is spin--spin ($\boldsymbol\sigma_1\cdot\boldsymbol\sigma_2$),
with \emph{no} $\hat{\mathbf r}$ dependence and no $1/r^2$ or $1/(\lambda r)$
terms -- it is $V_9,V_{10}$ (monopole--dipole, a single spin dotted into
$\hat{\mathbf r}$) that has that radial structure, not $V_2$. Second,
$V_{11}$ is $(\boldsymbol\sigma_1\times\boldsymbol\sigma_2)\cdot\hat{\mathbf r}$,
a cross product of the \emph{two} particles' spins with $\hat{\mathbf r}$,
not an $L\cdot S$ orbital spin-orbit coupling, while $V_{12},V_{13}$ are the
$\pm$ pair $(\boldsymbol\sigma_1\pm\boldsymbol\sigma_2)\cdot\mathbf v$.
Third, $V_4,V_5$, $V_9,V_{10}$, and $V_{12},V_{13}$ are each a $\pm$ pair of
the \emph{same} structure (sum vs.\ difference of the two particles' spins),
not independent structures.

\subsection{Classification under Discrete Symmetries}

Table~\ref{tab:dm-classification} groups the sixteen potentials by their
spin-structure class, which is the organising principle Dobrescu and Mocioiu
use to enumerate them exhaustively and is the property that determines
which SME coefficient(s) can generate a given $V_i$ in Chapter~3.

\begin{table}[h]
\centering
\begin{tabular}{@{}ll@{}}
\toprule
Structural class & Potentials \\
\midrule
Monopole--monopole (spin-independent) & $V_1$ \\
Monopole--dipole & $V_9, V_{10}$ \\
Dipole--dipole (spin--spin) & $V_2, V_3$ \\
Spin--velocity & $V_4, V_5, V_6, V_7$ \\
Velocity--velocity & $V_8$ \\
Spin cross-product & $V_{11}$ \\
Single-spin--velocity & $V_{12}, V_{13}, V_{14}$ \\
Higher-order spin--velocity & $V_{15}, V_{16}$ \\
\bottomrule
\end{tabular}
\caption{Structural classification of the Dobrescu--Mocioiu potentials.}
\label{tab:dm-classification}
\end{table}

\section{The Foldy--Wouthuysen Transformation}
\label{sec:fw}

\subsection{Purpose and General Method}

The FW transformation \cite{Foldy1950} is the standard technique for
systematically extracting the non-relativistic limit of a relativistic
Dirac Hamiltonian, order by order in $1/m$, while keeping every step exact
and unitary (so that no physical content is discarded, only reorganised by
inverse powers of the fermion mass). It is the bridge this thesis uses to
convert each SME coefficient's fully relativistic Lagrangian term into an
explicit non-relativistic Hamiltonian that can be matched onto the
Dobrescu--Mocioiu basis of Section~\ref{sec:dm}.

The method decouples the upper (particle) and lower (antiparticle)
two-component blocks of the four-component Dirac spinor by finding a
sequence of unitary transformations that eliminate, order by order, the
`odd' operators that mix the two blocks. Any perturbation $\varepsilon$ added
to the free Dirac Hamiltonian can be split into a part $\varepsilon_{\text{even}}$
that commutes with $\beta=\gamma^0$ (and is therefore already block-diagonal)
and a part $\varepsilon_{\text{odd}}$ that anticommutes with $\beta$ (and
therefore mixes the blocks), so that
\begin{equation}
  H = \beta m + \varepsilon_{\text{even}} + \varepsilon_{\text{odd}}.
\end{equation}
The FW generator chosen to eliminate $\varepsilon_{\text{odd}}$ at leading
order is
\begin{equation}
  U_{FW} = \exp\!\left(\frac{\beta\varepsilon_{\text{odd}}}{2m}\right),
\end{equation}
under which the Hamiltonian transforms as
\begin{equation}
  H' = U_{FW}^\dagger\, H\, U_{FW} - iU_{FW}^\dagger\,\partial_t U_{FW}.
\end{equation}

\subsection{General Order-by-Order Expansion}

Expanding the transformed Hamiltonian to order $1/m$ gives the master
expansion formula used throughout Chapter~3,
\begin{equation}
  H' \approx \beta m + \varepsilon_{\text{even}}
  + \frac{\beta\varepsilon_{\text{odd}}^2}{2m}
  - \frac{[\varepsilon_{\text{odd}},\varepsilon_{\text{even}}]}{4m}
  + O(m^{-2}). \label{eq:FWexpansion-ch2}
\end{equation}
Three features of Eq.~\eqref{eq:FWexpansion-ch2} govern every derivation in
Chapter~3. First, any even piece of the original perturbation,
$\varepsilon_{\text{even}}$, survives into the non-relativistic Hamiltonian
completely unchanged at order $m^0$ -- no FW iteration is required for it.
Second, an odd piece $\varepsilon_{\text{odd}}$ does not contribute at order
$m^0$ at all; its leading contribution appears only at order $1/m$, through
the $\beta\varepsilon_{\text{odd}}^2/2m$ term, which is why coefficients such
as $b_0$ and $H_{0i}$ that couple through odd operators are systematically
subleading (order $m^{-1}$) relative to coefficients such as $b_i$ and
$H_{ij}$ that couple through even operators (order $m^0$). Third, projecting
the block-diagonal result onto the upper $2\times2$ block (via $\beta\to+1$)
gives the single-particle, matter-sector non-relativistic Hamiltonian
directly; the antiparticle sector requires the separate charge-conjugation
argument developed in Chapter~3, Section~3.3, since reading off the lower
block of a single first-quantised Hamiltonian is not equivalent to
genuinely charge-conjugating the underlying field-theoretic coupling.

\subsection{Role in This Thesis}

Chapter~3 applies exactly this machinery -- unmodified -- to three SME
coefficients in turn: $b_\mu$ (Section 3.2), whose spatial component
$b_i$ is purely even and therefore contributes at leading order $m^0$;
$H_{\mu\nu}$ (Section 3.3), whose magnetic-like component $H_{ij}$ is
likewise even at leading order while the electric-like component $H_{0i}$ is
odd and therefore subleading; and $d_{\mu\nu}$ (Section 3.4), whose
mass-enhanced component $d_{i0}$ is even and agrees exactly with the
literature, while the momentum-dependent components $d_{ij}$ and $d_{00}$
are correct in operator structure but carry an overall sign this thesis's
equation-of-motion-level treatment does not yet resolve. In every
case, the resulting non-relativistic Hamiltonian is matched onto the
Dobrescu--Mocioiu catalogue of Section~\ref{sec:dm-catalogue} to identify
which $V_i$ the SME coefficient generates in a two-body interaction.

\section{Chapter Summary}

This chapter established the general theoretical machinery this thesis
relies on throughout: the minimal SME fermion Lagrangian and a classification
of its coefficients by CPT parity (Section~\ref{sec:sme}); the complete,
closed-form Dobrescu--Mocioiu catalogue of sixteen non-relativistic exotic
potentials (Section~\ref{sec:dm}); and the general Foldy--Wouthuysen
expansion used to move between the two (Section~\ref{sec:fw}). Chapter~3
applies this machinery to derive the explicit SME$\to$DM mapping for
$b_\mu$, $H_{\mu\nu}$, and $d_{\mu\nu}$, and Chapter~4 uses the resulting
matching table to interpret the compiled experimental constraint database.

\ifSubfilesClassLoaded{
\begin{thebibliography}{99}

\bibitem{Colladay1997}
D.~Colladay and V.~A. Kostelek\'{y},
\textit{CPT violation and the standard model},
Phys.\ Rev.\ D \textbf{55}, 6760 (1997).

\bibitem{Colladay1998}
D.~Colladay and V.~A. Kostelek\'{y},
\textit{Lorentz-violating extension of the standard model},
Phys.\ Rev.\ D \textbf{58}, 116002 (1998).

\bibitem{Dobrescu2006}
B.~A. Dobrescu and I.~Mocioiu,
\textit{Spin-dependent macroscopic forces from new particle exchange},
J.\ High Energy Phys.\ \textbf{11}, 005 (2006).

\bibitem{KosteleckyLane1999}
V.~A. Kostelek\'{y} and C.~D. Lane,
\textit{Nonrelativistic quantum Hamiltonian for Lorentz violation},
Phys.\ Rev.\ D \textbf{60}, 116010 (1999).

\bibitem{KosteleckyMewes2001}
V.~A. Kostelek\'{y} and M.~Mewes,
\textit{CPT violation and the standard model},
Phys.\ Rev.\ D \textbf{66}, 056005 (2001).

\bibitem{Foldy1950}
L.~L. Foldy and S.~A. Wouthuysen,
\textit{On the Dirac theory of spin-$\frac{1}{2}$ particles and its
non-relativistic limit},
Phys.\ Rev.\ \textbf{78}, 29 (1950).

\bibitem{Cong2025}
L.~Cong \textit{et al.},
\textit{Spin-dependent exotic interactions},
Rev.\ Mod.\ Phys.\ \textbf{97}, 025005 (2025).

\bibitem{Kostelecky2011}
V.~A. Kostelek\'{y} and N.~Russell,
\textit{Data tables for Lorentz and CPT violation},
Rev.\ Mod.\ Phys.\ \textbf{83}, 11 (2011), updated annually.

\end{thebibliography}
}{}

\end{document}
